\documentclass[11pt]{article}

\usepackage[margin=0.72in]{geometry}
\usepackage{amsmath,amssymb,amsfonts}
\usepackage{array}
\usepackage{booktabs}
\usepackage{enumitem}
\usepackage{tcolorbox}
\usepackage{xcolor}
\usepackage{fancyhdr}
\usepackage{lastpage}
\usepackage{multicol}
\usepackage{mathtools}

% =========================
% Novel Prep Colors
% =========================
\definecolor{novelPurple}{RGB}{98,54,255}
\definecolor{novelPurpleDark}{RGB}{76,42,204}
\definecolor{novelLight}{RGB}{242,238,255}
\definecolor{novelSoft}{RGB}{248,246,255}
\definecolor{darkGray}{RGB}{45,45,45}

% =========================
% Header / Footer
% =========================
\pagestyle{fancy}
\fancyhf{}
\lhead{\textcolor{novelPurpleDark}{\textbf{Novel Prep Learning Center}}}
\rhead{\textcolor{novelPurpleDark}{\textbf{AP Calculus AB \& BC Final Review}}}
\cfoot{\textcolor{novelPurpleDark}{\thepage\ of \pageref{LastPage}}}
\setlength{\headheight}{14pt}

\setlength{\parindent}{0pt}
\setlength{\parskip}{0.42em}
\renewcommand{\arraystretch}{1.18}

% =========================
% Boxes
% =========================
\newtcolorbox{unitbox}[1]{
  colback=novelLight,
  colframe=novelPurpleDark,
  title=\textbf{#1},
  fonttitle=\bfseries,
  coltitle=white,
  colbacktitle=novelPurpleDark,
  boxrule=0.85pt,
  arc=2pt,
  left=8pt,right=8pt,top=7pt,bottom=7pt
}

\newtcolorbox{formulabox}[1]{
  colback=white,
  colframe=novelPurple,
  title=\textbf{#1},
  fonttitle=\bfseries,
  coltitle=white,
  colbacktitle=novelPurple,
  boxrule=0.85pt,
  arc=1.5pt,
  left=8pt,right=8pt,top=7pt,bottom=7pt
}

\newtcolorbox{warningbox}[1]{
  colback=novelSoft,
  colframe=novelPurpleDark,
  title=\textbf{#1},
  fonttitle=\bfseries,
  coltitle=white,
  colbacktitle=novelPurpleDark,
  boxrule=0.85pt,
  arc=1.5pt,
  left=8pt,right=8pt,top=7pt,bottom=7pt
}

\newcommand{\sectiontitle}[1]{
\vspace{0.8em}
{\Large\textcolor{novelPurpleDark}{\textbf{#1}}}
\vspace{0.35em}
\hrule
\vspace{0.8em}
}

\newcommand{\subtopic}[1]{
\vspace{0.45em}
\textcolor{novelPurpleDark}{\textbf{#1}}
}

\begin{document}

% =====================================================
% COVER
% =====================================================

\begin{center}
    {\Huge \textbf{\textcolor{novelPurpleDark}{AP Calculus AB \& BC}}}\\[0.25em]
    {\LARGE \textbf{Complete Knowledge \& Formula Summary}}\\[0.5em]
    {\large Novel Prep Learning Center}\\[0.4em]
    {\normalsize Optimized AB / BC Separated Review Sheet}
\end{center}

\vspace{0.6em}

\begin{unitbox}{Purpose of This Summary}
This handout is designed for students who need a fast but complete AP Calculus review.  
Instead of rereading every unit from the beginning, students can use this document to rebuild the full knowledge structure, memorize core formulas, and identify missing topics quickly.

\medskip

\textbf{Important separation:}
\begin{itemize}[leftmargin=1.5em]
\item \textbf{AP Calculus AB Core:} Units 1--8 main concepts.
\item \textbf{AP Calculus BC Core:} all AB topics plus BC-only extensions, Unit 9, and Unit 10.
\item \textbf{BC-only topics placed separately:} logistic differential equations, Euler's Method, advanced integration techniques, parametric, polar, vector-valued functions, and infinite series.
\end{itemize}
\end{unitbox}

\newpage

% =====================================================
% AB SECTION
% =====================================================

\sectiontitle{Part I: AP Calculus AB Core Summary}

% =====================================================
\begin{unitbox}{AB Unit 1: Limits and Continuity}
\textbf{Big Idea:} A limit describes the behavior of a function as \(x\) approaches a value, not necessarily the value of the function at that point.
\end{unitbox}

\subtopic{1. Limit Notation}

\[
\lim_{x\to a} f(x)=L
\]

This means that as \(x\) gets closer to \(a\), the function values \(f(x)\) get closer to \(L\).

\subtopic{2. One-Sided Limits}

\[
\lim_{x\to a^-}f(x)
\qquad
\lim_{x\to a^+}f(x)
\]

A two-sided limit exists only if

\[
\lim_{x\to a^-}f(x)=\lim_{x\to a^+}f(x)
\]

\subtopic{3. Continuity at a Point}

A function \(f\) is continuous at \(x=a\) if:

\[
f(a)\text{ exists}
\]

\[
\lim_{x\to a}f(x)\text{ exists}
\]

\[
\lim_{x\to a}f(x)=f(a)
\]

\subtopic{4. Types of Discontinuities}

\begin{itemize}[leftmargin=1.5em]
\item \textbf{Removable discontinuity:} hole in the graph.
\item \textbf{Jump discontinuity:} left and right limits are different.
\item \textbf{Infinite discontinuity:} function approaches \(\pm\infty\).
\end{itemize}

\subtopic{5. Algebraic Limit Techniques}

\begin{itemize}[leftmargin=1.5em]
\item Direct substitution
\item Factoring and canceling
\item Rationalizing
\item Combining fractions
\item Squeeze Theorem
\end{itemize}

\begin{formulabox}{Squeeze Theorem}
If
\[
g(x)\le f(x)\le h(x)
\]
near \(x=a\), and

\[
\lim_{x\to a}g(x)=\lim_{x\to a}h(x)=L,
\]

then

\[
\lim_{x\to a}f(x)=L.
\]
\end{formulabox}

\subtopic{6. Limits at Infinity and Horizontal Asymptotes}

For rational functions:

\[
\lim_{x\to\infty}\frac{P(x)}{Q(x)}
\]

compare the highest powers.

\[
\deg P<\deg Q \Rightarrow y=0
\]

\[
\deg P=\deg Q \Rightarrow y=\frac{\text{leading coefficient of }P}{\text{leading coefficient of }Q}
\]

\[
\deg P>\deg Q \Rightarrow \text{no horizontal asymptote}
\]

\newpage

% =====================================================
\begin{unitbox}{AB Unit 2: Derivative Definition and Basic Rules}
\textbf{Big Idea:} The derivative measures instantaneous rate of change and slope of the tangent line.
\end{unitbox}

\subtopic{1. Derivative Definition}

\[
f'(x)=\lim_{h\to0}\frac{f(x+h)-f(x)}{h}
\]

Alternative form:

\[
f'(a)=\lim_{x\to a}\frac{f(x)-f(a)}{x-a}
\]

\subtopic{2. Tangent Line Equation}

If a function passes through \((a,f(a))\) and has slope \(f'(a)\), then

\[
y-f(a)=f'(a)(x-a)
\]

\subtopic{3. Basic Derivative Rules}

\[
\frac{d}{dx}[c]=0
\]

\[
\frac{d}{dx}[x^n]=nx^{n-1}
\]

\[
\frac{d}{dx}[cf(x)]=cf'(x)
\]

\[
\frac{d}{dx}[f(x)\pm g(x)]=f'(x)\pm g'(x)
\]

\subtopic{4. Trigonometric Derivatives}

\[
\frac{d}{dx}(\sin x)=\cos x
\]

\[
\frac{d}{dx}(\cos x)=-\sin x
\]

\[
\frac{d}{dx}(\tan x)=\sec^2 x
\]

\[
\frac{d}{dx}(\sec x)=\sec x\tan x
\]

\[
\frac{d}{dx}(\csc x)=-\csc x\cot x
\]

\[
\frac{d}{dx}(\cot x)=-\csc^2 x
\]

\subtopic{5. Product and Quotient Rules}

\[
(uv)'=u'v+uv'
\]

\[
\left(\frac{u}{v}\right)'=\frac{u'v-uv'}{v^2}
\]

\begin{warningbox}{Common Mistake}
The product rule is not
\[
(uv)'=u'v'.
\]
Always use
\[
u'v+uv'.
\]
\end{warningbox}

\newpage

% =====================================================
\begin{unitbox}{AB Unit 3: Composite, Implicit, and Inverse Function Derivatives}
\textbf{Big Idea:} Unit 3 handles functions with layers, hidden variables, and inverse relationships.
\end{unitbox}

\subtopic{1. Chain Rule}

\[
\frac{d}{dx}f(g(x))=f'(g(x))g'(x)
\]

Common forms:

\[
\frac{d}{dx}(u^n)=nu^{n-1}u'
\]

\[
\frac{d}{dx}(\sin u)=\cos u\cdot u'
\]

\[
\frac{d}{dx}(e^u)=e^u u'
\]

\[
\frac{d}{dx}(\ln u)=\frac{u'}{u}
\]

\subtopic{2. Exponential and Logarithmic Derivatives}

\[
\frac{d}{dx}(e^x)=e^x
\]

\[
\frac{d}{dx}(a^x)=a^x\ln a
\]

\[
\frac{d}{dx}(\ln x)=\frac1x
\]

\[
\frac{d}{dx}(\log_a x)=\frac1{x\ln a}
\]

\subtopic{3. Implicit Differentiation}

When \(y\) is not solved explicitly as a function of \(x\), differentiate both sides with respect to \(x\).

Important rule:

\[
\frac{d}{dx}(y^n)=ny^{n-1}\frac{dy}{dx}
\]

Example structure:

\[
x^2+y^2=25
\]

\[
2x+2y\frac{dy}{dx}=0
\]

\[
\frac{dy}{dx}=-\frac{x}{y}
\]

\subtopic{4. Derivative of an Inverse Function}

\[
(f^{-1})'(a)=\frac{1}{f'(f^{-1}(a))}
\]

Equivalent idea:

If \(f(b)=a\), then

\[
(f^{-1})'(a)=\frac{1}{f'(b)}
\]

\begin{warningbox}{Common Mistake}
For inverse derivative problems, students often plug into the wrong function.  
Always find the original \(x\)-value first.
\end{warningbox}

\newpage

% =====================================================
\begin{unitbox}{AB Unit 4: Contextual Applications of Differentiation}
\textbf{Big Idea:} Derivatives describe rates, motion, approximation, and relationships between changing quantities.
\end{unitbox}

\subtopic{1. Units of the Derivative}

If \(y=f(x)\), then

\[
f'(x)=\frac{\text{units of }y}{\text{units of }x}
\]

Example: If position is in feet and time is in seconds, velocity is measured in feet per second.

\subtopic{2. Motion}

If \(s(t)\) is position, then

\[
v(t)=s'(t)
\]

\[
a(t)=v'(t)=s''(t)
\]

\[
\text{speed}=|v(t)|
\]

\subtopic{3. Increasing and Decreasing Motion}

\[
v(t)>0 \Rightarrow \text{moving right/up}
\]

\[
v(t)<0 \Rightarrow \text{moving left/down}
\]

\[
a(t)>0 \Rightarrow v(t)\text{ increasing}
\]

\[
a(t)<0 \Rightarrow v(t)\text{ decreasing}
\]

\subtopic{4. Speeding Up and Slowing Down}

\[
v(t)\text{ and }a(t)\text{ same sign} \Rightarrow \text{speeding up}
\]

\[
v(t)\text{ and }a(t)\text{ opposite signs} \Rightarrow \text{slowing down}
\]

\subtopic{5. Related Rates}

General process:

\begin{enumerate}[leftmargin=1.5em]
\item Draw a diagram.
\item Write an equation relating the variables.
\item Differentiate both sides with respect to time.
\item Substitute known values after differentiating.
\item Solve for the requested rate.
\end{enumerate}

\subtopic{6. Linear Approximation}

\[
L(x)=f(a)+f'(a)(x-a)
\]

Differential form:

\[
dy=f'(x)dx
\]

\subtopic{7. L'Hospital's Rule}

For indeterminate forms

\[
\frac00
\qquad\text{or}\qquad
\frac{\infty}{\infty},
\]

if conditions are satisfied, then

\[
\lim_{x\to a}\frac{f(x)}{g(x)}
=
\lim_{x\to a}\frac{f'(x)}{g'(x)}
\]

\begin{warningbox}{Common Mistake}
Do not use L'Hospital's Rule unless the expression is actually an indeterminate form.
\end{warningbox}

\newpage

% =====================================================
\begin{unitbox}{AB Unit 5: Analytical Applications of Differentiation}
\textbf{Big Idea:} Derivatives describe the shape, behavior, and optimization of functions.
\end{unitbox}

\subtopic{1. Critical Numbers}

Critical numbers occur where

\[
f'(x)=0
\]

or

\[
f'(x)\text{ is undefined}
\]

provided \(x\) is in the domain of \(f\).

\subtopic{2. Extreme Value Theorem}

If \(f\) is continuous on a closed interval \([a,b]\), then \(f\) has an absolute maximum and absolute minimum on \([a,b]\).

\subtopic{3. Rolle's Theorem}

If \(f\) is continuous on \([a,b]\), differentiable on \((a,b)\), and

\[
f(a)=f(b),
\]

then there exists \(c\in(a,b)\) such that

\[
f'(c)=0.
\]

\subtopic{4. Mean Value Theorem}

If \(f\) is continuous on \([a,b]\) and differentiable on \((a,b)\), then there exists \(c\in(a,b)\) such that

\[
f'(c)=\frac{f(b)-f(a)}{b-a}.
\]

\subtopic{5. Increasing and Decreasing}

\[
f'(x)>0 \Rightarrow f\text{ increasing}
\]

\[
f'(x)<0 \Rightarrow f\text{ decreasing}
\]

\subtopic{6. First Derivative Test}

\[
f':+\to - \Rightarrow \text{relative maximum}
\]

\[
f':-\to + \Rightarrow \text{relative minimum}
\]

\subtopic{7. Concavity}

\[
f''(x)>0 \Rightarrow \text{concave up}
\]

\[
f''(x)<0 \Rightarrow \text{concave down}
\]

\subtopic{8. Inflection Points}

A point of inflection occurs where concavity changes.

Usually check candidates where

\[
f''(x)=0
\]

or

\[
f''(x)\text{ is undefined}.
\]

\subtopic{9. Second Derivative Test}

If \(f'(c)=0\), then:

\[
f''(c)>0 \Rightarrow \text{relative minimum}
\]

\[
f''(c)<0 \Rightarrow \text{relative maximum}
\]

\[
f''(c)=0 \Rightarrow \text{test inconclusive}
\]

\subtopic{10. Optimization}

General process:

\begin{enumerate}[leftmargin=1.5em]
\item Define variables.
\item Write the quantity to maximize or minimize.
\item Use constraints to write one-variable function.
\item Find critical points.
\item Check endpoints if interval is closed.
\item Interpret final answer in context.
\end{enumerate}

\newpage

% =====================================================
\begin{unitbox}{AB Unit 6: Integration and Accumulation}
\textbf{Big Idea:} Integration turns local rate of change into total accumulated change.
\end{unitbox}

\subtopic{1. Riemann Sums}

Approximate area:

\[
\sum_{i=1}^n f(x_i^*)\Delta x
\]

where

\[
\Delta x=\frac{b-a}{n}
\]

\subtopic{2. Definite Integral}

\[
\int_a^b f(x)\,dx
\]

represents signed area or accumulated change.

\subtopic{3. Fundamental Theorem of Calculus Part 1}

If

\[
F(x)=\int_a^x f(t)\,dt,
\]

then

\[
F'(x)=f(x).
\]

More generally:

\[
\frac{d}{dx}\int_a^{g(x)} f(t)\,dt=f(g(x))g'(x)
\]

\[
\frac{d}{dx}\int_{g(x)}^{h(x)} f(t)\,dt=f(h(x))h'(x)-f(g(x))g'(x)
\]

\subtopic{4. Fundamental Theorem of Calculus Part 2}

If \(F'(x)=f(x)\), then

\[
\int_a^b f(x)\,dx=F(b)-F(a)
\]

\subtopic{5. Basic Antiderivatives}

\[
\int x^n\,dx=\frac{x^{n+1}}{n+1}+C,\qquad n\neq -1
\]

\[
\int \frac1x\,dx=\ln|x|+C
\]

\[
\int e^x\,dx=e^x+C
\]

\[
\int a^x\,dx=\frac{a^x}{\ln a}+C
\]

\[
\int \cos x\,dx=\sin x+C
\]

\[
\int \sin x\,dx=-\cos x+C
\]

\[
\int \sec^2x\,dx=\tan x+C
\]

\subtopic{6. \(u\)-Substitution}

If \(u=g(x)\), then \(du=g'(x)\,dx\).

\[
\int f(g(x))g'(x)\,dx=\int f(u)\,du
\]

\begin{warningbox}{AB / BC Separation}
In this optimized version, AB Unit 6 includes basic antiderivatives and \(u\)-substitution.  
The following advanced integration techniques are moved to BC-only:
\[
\text{long division, completing the square, integration by parts, partial fractions, improper integrals.}
\]
\end{warningbox}

\newpage

% =====================================================
\begin{unitbox}{AB Unit 7: Differential Equations}
\textbf{Big Idea:} A differential equation tells how a function changes.
\end{unitbox}

\subtopic{1. Differential Equations}

A differential equation involves a function and its derivative.

Example:

\[
\frac{dy}{dx}=ky
\]

means the rate of change of \(y\) is proportional to \(y\).

\subtopic{2. Verifying Solutions}

To verify a solution:

\begin{enumerate}[leftmargin=1.5em]
\item Differentiate the proposed function.
\item Substitute into the differential equation.
\item Check the initial condition if given.
\end{enumerate}

\subtopic{3. Slope Fields}

A slope field shows small tangent segments representing

\[
\frac{dy}{dx}
\]

at many points.

\subtopic{4. Separation of Variables}

If the equation can be written as

\[
\frac{dy}{dx}=g(x)h(y),
\]

then separate:

\[
\frac{1}{h(y)}\,dy=g(x)\,dx
\]

and integrate both sides.

\subtopic{5. Exponential Growth and Decay}

\[
\frac{dy}{dt}=ky
\]

has solution

\[
y=Ce^{kt}
\]

If \(k>0\), growth.  
If \(k<0\), decay.

\begin{warningbox}{BC-only Topics Removed from AB Unit 7}
Euler's Method and logistic differential equations are listed in the BC-only section below.
\end{warningbox}

\newpage

% =====================================================
\begin{unitbox}{AB Unit 8: Applications of Integration}
\textbf{Big Idea:} Definite integrals measure average value, accumulation, area, and volume.
\end{unitbox}

\subtopic{1. Average Value of a Function}

\[
f_{\text{avg}}=\frac{1}{b-a}\int_a^b f(x)\,dx
\]

\subtopic{2. Net Change}

If \(r(t)\) is a rate, then

\[
\int_a^b r(t)\,dt
\]

gives net change.

\subtopic{3. Motion with Integrals}

\[
\text{displacement}=\int_a^b v(t)\,dt
\]

\[
\text{total distance}=\int_a^b |v(t)|\,dt
\]

\[
v(t)=v(a)+\int_a^t a(s)\,ds
\]

\[
s(t)=s(a)+\int_a^t v(s)\,ds
\]

\subtopic{4. Area Between Curves}

With respect to \(x\):

\[
A=\int_a^b [\text{top}-\text{bottom}]\,dx
\]

With respect to \(y\):

\[
A=\int_c^d [\text{right}-\text{left}]\,dy
\]

\subtopic{5. Volumes by Cross Sections}

\[
V=\int_a^b A(x)\,dx
\]

Common cross-sectional areas:

\[
\text{square: }A=s^2
\]

\[
\text{semicircle: }A=\frac{\pi}{8}s^2
\]

\[
\text{equilateral triangle: }A=\frac{\sqrt3}{4}s^2
\]

\subtopic{6. Disk and Washer Method}

Disk:

\[
V=\pi\int_a^b R(x)^2\,dx
\]

Washer:

\[
V=\pi\int_a^b \left(R(x)^2-r(x)^2\right)\,dx
\]

\begin{warningbox}{Common Mistake}
For washer method, do not subtract radii first.  
The correct structure is

\[
\pi\int (R^2-r^2)\,dx
\]

not

\[
\pi\int (R-r)^2\,dx.
\]
\end{warningbox}

\newpage

% =====================================================
% BC SECTION
% =====================================================

\sectiontitle{Part II: AP Calculus BC-Only Summary}

\begin{unitbox}{BC Overview}
AP Calculus BC includes all AB content plus additional topics involving advanced integration, differential equation approximation, parametric / polar / vector-valued calculus, and infinite series.

\medskip

The following sections are \textbf{BC-only} in this optimized review.
\end{unitbox}

% =====================================================
\begin{unitbox}{BC Extension A: Advanced Integration Techniques}
\textbf{Big Idea:} BC requires students to select integration techniques beyond basic \(u\)-substitution.
\end{unitbox}

\subtopic{1. Polynomial Long Division}

Use long division when the degree of the numerator is greater than or equal to the degree of the denominator.

Example structure:

\[
\int \frac{x^2+3x+5}{x+1}\,dx
\]

First divide:

\[
\frac{x^2+3x+5}{x+1}=x+2+\frac{3}{x+1}
\]

Then integrate:

\[
\int \left(x+2+\frac{3}{x+1}\right)\,dx
\]

\[
=\frac{x^2}{2}+2x+3\ln|x+1|+C
\]

\subtopic{2. Completing the Square}

Use completing the square when the denominator contains an irreducible quadratic.

Example:

\[
x^2+4x+8=(x+2)^2+4
\]

Useful forms:

\[
\int \frac{1}{u^2+a^2}\,du=\frac1a\tan^{-1}\left(\frac{u}{a}\right)+C
\]

\[
\int \frac{1}{\sqrt{a^2-u^2}}\,du=\sin^{-1}\left(\frac{u}{a}\right)+C
\]

\subtopic{3. Integration by Parts}

\[
\int u\,dv=uv-\int v\,du
\]

Use LIATE to choose \(u\):

\[
\text{Logarithmic} \to \text{Inverse trig} \to \text{Algebraic} \to \text{Trig} \to \text{Exponential}
\]

Common examples:

\[
\int x e^x\,dx
\]

\[
\int x\sin x\,dx
\]

\[
\int \ln x\,dx
\]

\subtopic{4. Partial Fractions}

Use partial fractions for rational functions after factoring the denominator.

\[
\frac{P(x)}{(x-a)(x-b)}
=
\frac{A}{x-a}+\frac{B}{x-b}
\]

Repeated linear factor:

\[
\frac{P(x)}{(x-a)^2}
=
\frac{A}{x-a}+\frac{B}{(x-a)^2}
\]

Irreducible quadratic factor:

\[
\frac{P(x)}{x^2+bx+c}
=
\frac{Ax+B}{x^2+bx+c}
\]

\subtopic{5. Improper Integrals}

Type 1: infinite interval.

\[
\int_a^\infty f(x)\,dx
=
\lim_{b\to\infty}\int_a^b f(x)\,dx
\]

Type 2: infinite discontinuity.

\[
\int_a^b f(x)\,dx
=
\lim_{t\to c^-}\int_a^t f(x)\,dx+
\lim_{t\to c^+}\int_t^b f(x)\,dx
\]

\begin{warningbox}{BC Integration Selection Guide}
\begin{itemize}[leftmargin=1.5em]
\item Product of different types of functions \(\Rightarrow\) integration by parts.
\item Rational function with high-degree numerator \(\Rightarrow\) long division first.
\item Rational function with factorable denominator \(\Rightarrow\) partial fractions.
\item Quadratic denominator that does not factor \(\Rightarrow\) complete the square.
\item Infinite bound or vertical asymptote \(\Rightarrow\) improper integral.
\end{itemize}
\end{warningbox}

\newpage

% =====================================================
\begin{unitbox}{BC Extension B: Euler's Method}
\textbf{Big Idea:} Euler's Method approximates a solution curve using tangent line steps.
\end{unitbox}

Given

\[
\frac{dy}{dx}=f(x,y)
\]

and an initial condition

\[
(x_0,y_0),
\]

Euler's Method uses

\[
y_{\text{new}}=y_{\text{old}}+(\text{slope})(\Delta x)
\]

where

\[
\text{slope}=f(x_{\text{old}},y_{\text{old}})
\]

So the update formula is

\[
y_{n+1}=y_n+f(x_n,y_n)\Delta x
\]

\[
x_{n+1}=x_n+\Delta x
\]

\begin{warningbox}{Euler's Method Common Mistake}
Students often use the original initial slope for every step.  
The slope must be recalculated at each new point.
\end{warningbox}

\newpage

% =====================================================
\begin{unitbox}{BC Extension C: Logistic Differential Equations}
\textbf{Big Idea:} Logistic growth models limited growth with a carrying capacity.
\end{unitbox}

\subtopic{1. Logistic Differential Equation}

\[
\frac{dy}{dt}=ky\left(1-\frac{y}{L}\right)
\]

where

\[
L=\text{carrying capacity}
\]

\[
k=\text{growth constant}
\]

\subtopic{2. Behavior}

\[
0<y<L \Rightarrow \frac{dy}{dt}>0
\]

The population increases.

\[
y>L \Rightarrow \frac{dy}{dt}<0
\]

The population decreases toward \(L\).

\subtopic{3. Maximum Growth Rate}

For logistic growth,

\[
\frac{dy}{dt}
\]

is maximized when

\[
y=\frac{L}{2}
\]

\subtopic{4. Logistic Solution Form}

A common solution form is

\[
y=\frac{L}{1+Ae^{-kt}}
\]

where \(A\) is determined by the initial condition.

\begin{warningbox}{Logistic Model Interpretation}
Logistic growth is not unlimited exponential growth.  
It begins similarly to exponential growth but eventually levels off near the carrying capacity.
\end{warningbox}

\newpage

% =====================================================
\begin{unitbox}{BC Unit 9: Parametric, Polar, and Vector-Valued Functions}
\textbf{Big Idea:} Curves can be described using parameters, vectors, or polar coordinates instead of only \(y=f(x)\).
\end{unitbox}

\subtopic{1. Parametric Equations}

A curve may be written as

\[
x=x(t),\qquad y=y(t)
\]

\subtopic{2. Parametric Derivative}

\[
\frac{dy}{dx}=\frac{dy/dt}{dx/dt}
\]

provided

\[
\frac{dx}{dt}\neq0.
\]

\subtopic{3. Second Derivative for Parametric Equations}

\[
\frac{d^2y}{dx^2}
=
\frac{\frac{d}{dt}\left(\frac{dy}{dx}\right)}{\frac{dx}{dt}}
\]

\subtopic{4. Parametric Arc Length}

\[
L=\int_a^b \sqrt{\left(\frac{dx}{dt}\right)^2+\left(\frac{dy}{dt}\right)^2}\,dt
\]

\subtopic{5. Vector-Valued Functions}

\[
\vec r(t)=\langle x(t),y(t)\rangle
\]

Velocity:

\[
\vec v(t)=\vec r'(t)
\]

Acceleration:

\[
\vec a(t)=\vec v'(t)=\vec r''(t)
\]

Speed:

\[
|\vec v(t)|=\sqrt{(x'(t))^2+(y'(t))^2}
\]

Total distance:

\[
\int_a^b |\vec v(t)|\,dt
\]

Displacement:

\[
\vec r(b)-\vec r(a)
\]

\subtopic{6. Polar Coordinates}

\[
x=r\cos\theta
\]

\[
y=r\sin\theta
\]

\[
r^2=x^2+y^2
\]

\[
\tan\theta=\frac{y}{x}
\]

\subtopic{7. Polar Slope}

If \(r=f(\theta)\), then

\[
\frac{dy}{dx}
=
\frac{\frac{dy}{d\theta}}{\frac{dx}{d\theta}}
\]

where

\[
x=r\cos\theta,\qquad y=r\sin\theta.
\]

So

\[
\frac{dx}{d\theta}=r'\cos\theta-r\sin\theta
\]

\[
\frac{dy}{d\theta}=r'\sin\theta+r\cos\theta
\]

Thus

\[
\frac{dy}{dx}
=
\frac{r'\sin\theta+r\cos\theta}{r'\cos\theta-r\sin\theta}
\]

\subtopic{8. Polar Area}

Area inside one polar curve:

\[
A=\frac12\int_\alpha^\beta r^2\,d\theta
\]

Area between two polar curves:

\[
A=\frac12\int_\alpha^\beta \left(r_{\text{outer}}^2-r_{\text{inner}}^2\right)d\theta
\]

\begin{warningbox}{Polar Area Common Mistake}
Do not use

\[
\int r\,d\theta.
\]

Polar area always uses

\[
\frac12\int r^2\,d\theta.
\]
\end{warningbox}

\newpage

% =====================================================
\begin{unitbox}{BC Unit 10: Infinite Sequences and Series}
\textbf{Big Idea:} Infinite series are studied through convergence, error control, and function representation.
\end{unitbox}

\subtopic{1. Sequence vs. Series}

A sequence is a list:

\[
\{a_n\}
\]

A series is a sum:

\[
\sum_{n=1}^{\infty}a_n
\]

Partial sum:

\[
S_n=a_1+a_2+\cdots+a_n
\]

\subtopic{2. Geometric Series}

\[
\sum_{n=0}^{\infty}ar^n
\]

Converges if

\[
|r|<1
\]

and the sum is

\[
S=\frac{a}{1-r}
\]

Diverges if

\[
|r|\ge1.
\]

\subtopic{3. nth-Term Test for Divergence}

If

\[
\lim_{n\to\infty}a_n\neq0
\]

or the limit does not exist, then

\[
\sum a_n
\]

diverges.

\begin{warningbox}{Important}
If

\[
\lim_{n\to\infty}a_n=0,
\]

the series may converge or diverge.  
The nth-term test cannot prove convergence.
\end{warningbox}

\subtopic{4. p-Series}

\[
\sum_{n=1}^{\infty}\frac1{n^p}
\]

Converges if

\[
p>1
\]

Diverges if

\[
p\le1
\]

\subtopic{5. Integral Test}

If \(f(x)\) is positive, continuous, and decreasing, and

\[
a_n=f(n),
\]

then

\[
\sum_{n=1}^{\infty}a_n
\]

and

\[
\int_1^\infty f(x)\,dx
\]

either both converge or both diverge.

\subtopic{6. Direct Comparison Test}

If

\[
0\le a_n\le b_n,
\]

then:

\[
\sum b_n \text{ converges} \Rightarrow \sum a_n \text{ converges}
\]

\[
\sum a_n \text{ diverges} \Rightarrow \sum b_n \text{ diverges}
\]

\subtopic{7. Limit Comparison Test}

For positive series:

\[
L=\lim_{n\to\infty}\frac{a_n}{b_n}
\]

If

\[
0<L<\infty,
\]

then

\[
\sum a_n
\]

and

\[
\sum b_n
\]

have the same behavior.

\subtopic{8. Alternating Series Test}

For

\[
\sum (-1)^n b_n
\]

or

\[
\sum (-1)^{n+1}b_n,
\]

the series converges if:

\[
b_n>0
\]

\[
b_n\text{ decreases}
\]

\[
\lim_{n\to\infty}b_n=0
\]

\subtopic{9. Alternating Series Error Bound}

If an alternating series satisfies the conditions of the Alternating Series Test, then

\[
|R_n|\le b_{n+1}
\]

\subtopic{10. Ratio Test}

\[
L=\lim_{n\to\infty}\left|\frac{a_{n+1}}{a_n}\right|
\]

\[
L<1 \Rightarrow \text{absolute convergence}
\]

\[
L>1 \Rightarrow \text{divergence}
\]

\[
L=1 \Rightarrow \text{inconclusive}
\]

\subtopic{11. Absolute vs. Conditional Convergence}

\[
\sum a_n \text{ absolutely converges if } \sum |a_n| \text{ converges.}
\]

\[
\sum a_n \text{ conditionally converges if } \sum a_n \text{ converges but } \sum |a_n| \text{ diverges.}
\]

\subtopic{12. Power Series}

\[
\sum_{n=0}^{\infty}c_n(x-a)^n
\]

Radius of convergence:

\[
R
\]

Interval of convergence:

\[
(a-R,a+R)
\]

Always check endpoints separately.

\subtopic{13. Taylor and Maclaurin Series}

Taylor series centered at \(x=a\):

\[
f(x)=\sum_{n=0}^{\infty}\frac{f^{(n)}(a)}{n!}(x-a)^n
\]

Maclaurin series centered at \(x=0\):

\[
f(x)=\sum_{n=0}^{\infty}\frac{f^{(n)}(0)}{n!}x^n
\]

\subtopic{14. Common Maclaurin Series}

\[
e^x=\sum_{n=0}^{\infty}\frac{x^n}{n!}
\]

\[
\sin x=\sum_{n=0}^{\infty}(-1)^n\frac{x^{2n+1}}{(2n+1)!}
\]

\[
\cos x=\sum_{n=0}^{\infty}(-1)^n\frac{x^{2n}}{(2n)!}
\]

\[
\frac1{1-x}=\sum_{n=0}^{\infty}x^n,\qquad |x|<1
\]

\[
\ln(1+x)=\sum_{n=1}^{\infty}(-1)^{n+1}\frac{x^n}{n},\qquad -1<x\le1
\]

\[
\tan^{-1}x=\sum_{n=0}^{\infty}(-1)^n\frac{x^{2n+1}}{2n+1},\qquad |x|\le1
\]

\subtopic{15. Lagrange Error Bound}

If \(P_n(x)\) is the Taylor polynomial for \(f(x)\), then

\[
|R_n(x)|\le \frac{M}{(n+1)!}|x-a|^{n+1}
\]

where \(M\) is an upper bound for

\[
|f^{(n+1)}(x)|
\]

on the interval between \(a\) and \(x\).

\newpage

% =====================================================
% COVERAGE CHECK
% =====================================================

\sectiontitle{Part III: AB / BC Coverage Check}

\begin{unitbox}{AB Topics Included}
\begin{multicols}{2}
\begin{itemize}[leftmargin=1.2em]
\item Limits
\item One-sided limits
\item Continuity
\item Discontinuities
\item Infinite limits
\item Limits at infinity
\item Derivative definition
\item Tangent lines
\item Basic derivative rules
\item Product rule
\item Quotient rule
\item Chain rule
\item Implicit differentiation
\item Inverse function derivatives
\item Exponential derivatives
\item Logarithmic derivatives
\item Related rates
\item Motion
\item Linear approximation
\item L'Hospital's Rule
\item Rolle's Theorem
\item Mean Value Theorem
\item Extreme Value Theorem
\item First derivative test
\item Second derivative test
\item Concavity
\item Optimization
\item Riemann sums
\item Definite integrals
\item FTC Part 1
\item FTC Part 2
\item Basic antiderivatives
\item \(u\)-substitution
\item Slope fields
\item Separable differential equations
\item Exponential growth / decay
\item Average value
\item Area between curves
\item Volume by cross sections
\item Disk / washer method
\item Displacement / total distance
\end{itemize}
\end{multicols}
\end{unitbox}

\begin{unitbox}{BC-Only Topics Included}
\begin{multicols}{2}
\begin{itemize}[leftmargin=1.2em]
\item Polynomial long division integration
\item Completing the square integrals
\item Integration by parts
\item Partial fractions
\item Improper integrals
\item Euler's Method
\item Logistic differential equations
\item Parametric equations
\item Parametric first derivative
\item Parametric second derivative
\item Parametric arc length
\item Vector-valued functions
\item Velocity vectors
\item Acceleration vectors
\item Speed
\item Displacement
\item Total distance
\item Polar coordinates
\item Polar slope
\item Polar area
\item Sequences
\item Infinite series
\item Geometric series
\item nth-term test
\item p-series
\item Integral test
\item Direct comparison
\item Limit comparison
\item Alternating series
\item Alternating series error bound
\item Ratio test
\item Absolute convergence
\item Conditional convergence
\item Power series
\item Radius of convergence
\item Interval of convergence
\item Taylor series
\item Maclaurin series
\item Lagrange error bound
\end{itemize}
\end{multicols}
\end{unitbox}

\begin{warningbox}{Remaining Topics Worth Adding Later}
This optimized summary covers the major AB and BC testable topics.  
For an even stronger final version, the following can be added as extra pages:

\begin{itemize}[leftmargin=1.5em]
\item A one-page derivative formula sheet.
\item A one-page integral formula sheet.
\item A series test decision flowchart.
\item A motion interpretation table.
\item A polar / parametric formula mini-sheet.
\item A common FRQ template page.
\item Calculator-active vs. non-calculator strategy page.
\end{itemize}
\end{warningbox}

\end{document}